% ENEM 2013-2019 — 6 questões MCQ — Função Exponencial e Função Logarítmica
% Fonte: listas "Função Exponencial" e "Função Logarítmica" do site projetoagathaedu.com.br
% Omitidas Exponencial: Q2,Q3,Q4,Q5,Q10,Q12 (imagens), Q7 (imagens nas alternativas), Q11 (duplicata lista14/Q12), Q13 (duplicata Q10)
% Omitidas Logarítmica: Q1,Q5 (imagem — tabela de classificações/efeitos), Q2 (imagem + duplicata Exp/Q3), Q3 (duplicata lista14/Q17), Q6,Q9 (fórmula em imagem), Q8 (alternativas com imagens)

---
tipo: Múltipla Escolha
dificuldade: Fácil
nivel: médio
disciplina: Matemática
assunto: Função Exponencial
gabarito: D
tags: [bactérias, epidemia, reprodução, duplicação]
fonte: concurso
concurso: ENEM
banca: INEP
ano: 2016
numero: 1
---
\question
O governo de uma cidade está preocupado com a possível epidemia de uma doença infectocontagiosa causada por bactéria. Para decidir que medidas tomar, deve calcular a velocidade de reprodução da bactéria. Em experiências laboratoriais de uma cultura bacteriana, inicialmente com 40 mil unidades, obteve-se a fórmula para a população:
\[ p(t) = 40 \cdot 2^{3t} \]
em que \( t \) é o tempo, em hora, e \( p(t) \) é a população, em milhares de bactérias.

Em relação à quantidade inicial de bactérias, após 20 min, a população será
\begin{choices}
\choice reduzida a um terço.
\choice reduzida à metade.
\choice reduzida a dois terços.
\correctchoice duplicada.
\choice triplicada.
\end{choices}


---
tipo: Múltipla Escolha
dificuldade: Média
nivel: médio
disciplina: Matemática
assunto: Função Exponencial
gabarito: C
tags: [bactérias, arsênico, crescimento exponencial, triplicar]
fonte: concurso
concurso: ENEM-PPL
banca: INEP
ano: 2019
numero: 6
---
\question
Em um laboratório, cientistas observaram o crescimento de uma população de bactérias submetida a uma dieta magra em fósforo, com generosas porções de arsênico. Descobriu-se que o número de bactérias dessa população, após \( t \) horas de observação, poderia ser modelado pela função exponencial \( N(t) = N_0 e^{kt} \), em que \( N_0 \) é o número de bactérias no instante do início da observação (\( t = 0 \)) e representa uma constante real maior que 1, e \( k \) é uma constante real positiva.

Sabe-se que, após uma hora de observação, o número de bactérias foi triplicado.

Cinco horas após o início da observação, o número de bactérias, em relação ao número inicial dessa cultura, foi
\begin{oneparchoices}
\choice \( 3N_0 \)
\choice \( 15N_0 \)
\correctchoice \( 243N_0 \)
\choice \( 360N_0 \)
\choice \( 729N_0 \)
\end{oneparchoices}

---
tipo: Múltipla Escolha
dificuldade: Fácil
nivel: médio
disciplina: Matemática
assunto: Função Exponencial
gabarito: E
tags: [salário, sindicato, piso salarial, reajuste]
fonte: concurso
concurso: ENEM-PPL
banca: INEP
ano: 2015
numero: 8
---
\question
O sindicato de trabalhadores de uma empresa sugere que o piso salarial da classe seja de R\$ 1{.}800,00, propondo um aumento percentual fixo por cada ano dedicado ao trabalho. A expressão que corresponde à proposta salarial \( s \), em função do tempo de serviço \( t \), em anos, é \( s(t) = 1{.}800 \cdot (1{,}03)^t \).

De acordo com a proposta do sindicato, o salário de um profissional dessa empresa com 2 anos de tempo de serviço será, em reais,
\begin{oneparchoices}
\choice 7{.}416,00.
\choice 3{.}819,24.
\choice 3{.}709,62.
\choice 3{.}708,00.
\correctchoice 1{.}909,62.
\end{oneparchoices}

---
tipo: Múltipla Escolha
dificuldade: Fácil
nivel: médio
disciplina: Matemática
assunto: Função Exponencial
gabarito: E
tags: [bactérias, bactericida, meia-vida, classificação de funções]
fonte: concurso
concurso: ENEM-PPL
banca: INEP
ano: 2013
numero: 9
---
\question
Em um experimento, uma cultura de bactérias tem sua população reduzida pela metade a cada hora, devido à ação de um agente bactericida.

Neste experimento, o número de bactérias em função do tempo pode ser modelado por uma função do tipo
\begin{choices}
\choice afim.
\choice seno.
\choice cosseno.
\choice logarítmica crescente.
\correctchoice exponencial.
\end{choices}

---
tipo: Múltipla Escolha
dificuldade: Difícil
nivel: médio
disciplina: Matemática
assunto: Função Logarítmica
gabarito: C
tags: [transistores, Lei de Moore, processador, computação]
fonte: concurso
concurso: ENEM
banca: INEP
ano: 2018
numero: 4
---
\question
Com o avanço em ciência da computação, estamos próximos do momento em que o número de transistores no processador de um computador pessoal será da mesma ordem de grandeza que o número de neurônios em um cérebro humano, que é da ordem de 100 bilhões.

Uma das grandezas determinantes para o desempenho de um processador é a densidade de transistores, que é o número de transistores por centímetro quadrado. Em 1986, uma empresa fabricava um processador contendo 100{.}000 transistores distribuídos em 0,25 cm² de área. Desde então, o número de transistores por centímetro quadrado que se pode colocar em um processador dobra a cada dois anos (Lei de Moore).

Considere 0,30 como aproximação para \( \log_{10} 2 \).

Em que ano a empresa atingiu ou atingirá a densidade de 100 bilhões de transistores?
\begin{oneparchoices}
\choice 1999
\choice 2002
\correctchoice 2022
\choice 2026
\choice 2146
\end{oneparchoices}

---
tipo: Múltipla Escolha
dificuldade: Difícil
nivel: médio
disciplina: Matemática
assunto: Função Logarítmica
gabarito: D
tags: [liga metálica, resfriamento, temperatura, logaritmo]
fonte: concurso
concurso: ENEM
banca: INEP
ano: 2016
numero: 7
---
\question
Uma liga metálica sai do forno a uma temperatura de 3{.}000 °C e diminui 1\% de sua temperatura a cada 30 min. Use 0,477 como aproximação para \( \log_{10}(3) \) e 1,041 como aproximação para \( \log_{10}(11) \).

O tempo decorrido, em hora, até que a liga atinja 30 °C é mais próximo de
\begin{oneparchoices}
\choice 22
\choice 50
\choice 100
\correctchoice 200
\choice 400
\end{oneparchoices}
